%-------------------------
% Resume in Latex
% Author : Yagiz Erdem
% Based off of: https://github.com/sb2nov/resume
% License : MIT
%------------------------

\documentclass[letterpaper,11pt]{article}

\usepackage{latexsym}
\usepackage[empty]{fullpage}
\usepackage{titlesec}
\usepackage{marvosym}
\usepackage[usenames,dvipsnames]{color}
\usepackage{verbatim}
\usepackage{enumitem}
\usepackage[hidelinks]{hyperref}
\usepackage{fancyhdr}
\usepackage[english]{babel}
\usepackage{tabularx}
\input{glyphtounicode}


%----------FONT OPTIONS----------
% sans-serif
% \usepackage[sfdefault]{FiraSans}
% \usepackage[sfdefault]{roboto}
% \usepackage[sfdefault]{noto-sans}
% \usepackage[default]{sourcesanspro}

% serif
% \usepackage{CormorantGaramond}
% \usepackage{charter}


\pagestyle{fancy}
\fancyhf{} % clear all header and footer fields
\fancyfoot{}
\renewcommand{\headrulewidth}{0pt}
\renewcommand{\footrulewidth}{0pt}

% Adjust margins
\addtolength{\oddsidemargin}{-0.5in}
\addtolength{\evensidemargin}{-0.5in}
\addtolength{\textwidth}{1in}
\addtolength{\topmargin}{-.5in}
\addtolength{\textheight}{1.0in}

\urlstyle{same}

\raggedbottom
\raggedright
\setlength{\tabcolsep}{0in}

% Sections formatting
\titleformat{\section}{
  \vspace{-4pt}\scshape\raggedright\large
}{}{0em}{}[\color{black}\titlerule \vspace{-5pt}]

% Ensure that generate pdf is machine readable/ATS parsable
\pdfgentounicode=1

%-------------------------
% Custom commands
\newcommand{\resumeItem}[1]{
  \item\small{
    {#1 \vspace{-2pt}}
  }
}

\newcommand{\resumeSubheading}[4]{
  \vspace{-2pt}\item
    \begin{tabular*}{0.97\textwidth}[t]{l@{\extracolsep{\fill}}r}
      \textbf{#1} & #2 \\
      \textit{\small#3} & \textit{\small #4} \\
    \end{tabular*}\vspace{-7pt}
}

\newcommand{\resumeSubSubheading}[2]{
    \item
    \begin{tabular*}{0.97\textwidth}{l@{\extracolsep{\fill}}r}
      \textit{\small#1} & \textit{\small #2} \\
    \end{tabular*}\vspace{-7pt}
}

\newcommand{\resumeProjectHeading}[2]{
    \item
    \begin{tabular*}{0.97\textwidth}{l@{\extracolsep{\fill}}r}
      \small#1 & #2 \\
    \end{tabular*}\vspace{-7pt}
}

\newcommand{\resumeSubItem}[1]{\resumeItem{#1}\vspace{-4pt}}

\renewcommand\labelitemii{$\vcenter{\hbox{\tiny$\bullet$}}$}

\newcommand{\resumeSubHeadingListStart}{\begin{itemize}[leftmargin=0.15in, label={}]}
\newcommand{\resumeSubHeadingListEnd}{\end{itemize}}
\newcommand{\resumeItemListStart}{\begin{itemize}}
\newcommand{\resumeItemListEnd}{\end{itemize}\vspace{-5pt}}

%-------------------------------------------
%%%%%%  RESUME STARTS HERE  %%%%%%%%%%%%%%%%%%%%%%%%%%%%


\begin{document}

\begin{center}
    \textbf{\Huge \scshape Yağız Erdem} \\ \vspace{1pt}
    \small PHONE-PLACEHOLDER $|$ \href{mailto:yagizerdem819@gmail.com}{\underline{yagizerdem819@gmail.com}} $|$
    \href{https://www.linkedin.com/in/erdemyagiz/}{\underline{linkedin.com/in/erdemyagiz}} $|$
    \href{https://github.com/yagizerdem}{\underline{github.com/yagizerdem}}
\end{center}

%-----------EDUCATION-----------
\section{Eğitim}
\resumeSubHeadingListStart

  \resumeSubheading
    {Dokuz Eylül Üniversitesi}{İzmir, Türkiye}
    {Bilgisayar Mühendisliği Lisans Programı}{2023 -- 2027}

  \resumeSubheading
    {Bahçeşehir Koleji}{İstanbul, Türkiye}
    {Lise Diploması}{2018 -- 2022}

\resumeSubHeadingListEnd

%-----------EXPERIENCE-----------
\section{Deneyim}
\resumeSubHeadingListStart

  \resumeSubheading
    {TypeScript Geliştirici Stajyeri}{Haz. 2026 -- Ağu. 2026}
    {Eliar Elektronik}{İstanbul, Türkiye}
    \resumeItemListStart
      \resumeItem{Büyük ölçekli bir monorepo ortamındaki TypeScript tabanlı uygulamalara üretim kodu katkısında bulundum; Vue.js, Nuxt.js, TypeScript ve Quasar kullanarak frontend özellikleri geliştirdim ve mevcut özelliklerin bakımını yaptım. Nuxt.js ekosistemi içerisinde backend API geliştirme ve uygulama mantığına katkıda bulundum, CI/CD iş akışları ile otomatik derleme ve dağıtım süreçlerinin geliştirilmesi ve bakımına destek oldum; hata ayıklama, özellik geliştirme ve mevcut üretim uygulamalarına yeni işlevlerin entegrasyonu üzerinde çalıştım.}
    \resumeItemListEnd

  \resumeSubheading
    {Yazılım Mühendisliği Stajyeri}{Oca. 2025 -- Şub. 2025}
    {Dokuz Eylül Üniversitesi}{İzmir, Türkiye}
    \resumeItemListStart
      \resumeItem{C\# ve .NET Core kullanarak komut satırı tabanlı bir Yer Radarı (GPR) veri işleme uygulaması olan GeoExtractor'ı geliştirdim. Ham MALÅ RD3/RD7 radar verilerini yapılandırılmış JSON formatına dönüştüren bir ayrıştırıcı geliştirdim, sinyal işleme algoritmalarını çalıştırmak için R Portable entegrasyonu gerçekleştirdim, işlenmiş radar çıktılarının Base64'ten PNG'ye otomatik dönüştürülmesini sağladım ve radar görüntüleri ile sinyal özelliklerini SQLite üzerinde sakladım. System.Diagnostics.Process kullanarak çoklu süreç yürütme mekanizmasını geliştirdim.}
    \resumeItemListEnd

  \resumeSubheading
    {Yazılım Mühendisliği Stajyeri}{Ağu. 2025 -- Eyl. 2025}
    {GameLab}{İstanbul, Türkiye}
    \resumeItemListStart
      \resumeItem{C++ kullanılarak geliştirilen Datuxa veri analizi uygulamasına katkıda bulundum; eksik, geçersiz ve tutarsız değerleri tespit edip işlemek için CSV veri temizleme işlevleri geliştirdim. Ayrıca masaüstü kullanıcı arayüzü geliştirme çalışmalarına katkıda bulundum ve mevcut arayüz bileşenlerinin kullanılabilirliğini iyileştirdim.}
    \resumeItemListEnd

\resumeSubHeadingListEnd

%-----------PROJECTS-----------
\section{Projeler}
\resumeSubHeadingListStart

  \resumeProjectHeading
    {\textbf{YSharp} $|$ \emph{Java, Yorumlayıcı, Recursive Descent Parser, AST}}{}

  \small{
    Dinamik tiplemeye ve C benzeri bir sözdizimine sahip yorumlanan bir betik dilini
    Java ile sıfırdan tasarladım ve geliştirdim. Özel bir lexer, recursive-descent parser,
    AST ve tree-walk evaluator dahil olmak üzere yorumlayıcının tüm çalışma hattını
    geliştirdim. Birinci sınıf fonksiyonlar, prototip tabanlı nesneler, diziler, hash map'ler,
    kontrol akışı ve hata yönetimi özelliklerini geliştirdim.
  }
  \vspace{6pt}

  \resumeProjectHeading
    {\textbf{Özel Regex Motoru} $|$ \emph{Java, Derleyici Tasarımı, NFA}}{}

  \small{
    Java'nın yerleşik regex implementasyonunu kullanmadan sıfırdan bir düzenli ifade
    motoru geliştirdim. Düzenli ifadeleri ayrıştırıp özel bir Belirlenimsiz Sonlu Otomat
    (NFA) yapısına derleyen bir derleyici hattı tasarladım. Eşleştirme işlemlerini
    geleneksel backtracking kullanmadan NFA simülasyonu üzerinden gerçekleştirdim.
  }
  \vspace{6pt}

  \resumeProjectHeading
    {\textbf{Mano Assembly Derleyicisi \& Simülatörü} $|$
    \emph{TypeScript, React, Bilgisayar Mimarisi}}{}

  \small{
    Bilgisayar mimarisi eğitimi için kapsamlı bir Mano Assembly geliştirme araç zinciri
    oluşturdum. Lexer, parser ve iki geçişli assembler bileşenlerini TypeScript ile
    sıfırdan geliştirdim. Komutların yürütülmesini ve işlemci durumundaki değişiklikleri
    görselleştirmek için mikro-adım seviyesinde çalışan bir CPU simülatörü geliştirdim.
  }

\resumeSubHeadingListEnd

%-----------PROGRAMMING SKILLS-----------
\section{Teknik Beceriler}
\begin{itemize}[leftmargin=0.15in, label={}]
  \small{\item{
    \textbf{Programlama Dilleri}{: Java, TypeScript, JavaScript, C, C++, C\#, Python, SQL, Bash} \\
    \textbf{Frameworkler \& Teknolojiler}{: Vue.js, Nuxt.js, React, Node.js, .Express.js, SQL, SQLite} \\
    \textbf{Geliştirici Araçları}{: Git, Docker, GitHub Actions, Gradle, Maven, pnpm, IntelliJ IDEA, VS Code} \\
    \textbf{Web \& Mobil}{: REST API'leri, Responsive Web Geliştirme, Mobil Uygulama Geliştirme, Kimlik Doğrulama,       Veritabanı Entegrasyonu, CI/CD}
  }}
\end{itemize}


%-------------------------------------------
\end{document}